start a tex file in here with all PO studies and PSPS studies with notes on findings from each 
- main epi studies and details written out of what they found 
- supporting studies: write out a sentence for each source but don't need that many specific findings all written out 

- make overleaf proj for notes from all studies i read and maybe sentences for papers. for auxilary studies, write the sentence right away so you have key info but dont have to keep going back to over and over. 


Next time: 
include years and study population and health dataset that was used. you want to be able to say whether you looked at the same locations/times
include the actual effect estimate for the main papers
reach goal: think about what is the sentence you want to write about this paper in my paper

CVD hospitalizations: lin2021immediate et al. null; Do et al. null; deng2022independent hosp positive; 
Resp hospitalizations: lin2021immediate et al. null; dominianni et al. positive; lin2024joint positive ED; deng2022independent hosp positive; 
COPD hospitalizations: zhang et al. positive; lin2024joint positive; 
Food/water borne hospitalizations: lin et al. null; deng2022independent hosp positive; 
injury hospitalizations: lin et al. null.; northrop et al. kids positive; 
Renal hospitalizations: dominianni et al. positive; 
mortality: dominianni et al. positive; skarha et al. positive; anderson2011impacts positive; 
CO poisoning: northrop et al. kids positive (northrop2024power)

NOTE: large vs. small power outages in dominianni 
TO DO: check whether each of these looks at ED or hospitalizations or both and amend this summary. 
TO DO: go look at the references for vivian's paper and check references to make sure im not missing anything. 

how do people measure power outages? synthesis of the synthesis: 
binary metric of power being out based on percent cutpoints of the percent of people being without power in the geo unit of analysis. -- alex, vivian, lin, lin, zhang, deng
anderson and dominniani just used the event of the 2003 outage and called everyone exposed. 
skarha used facility level power outage information 


- topic sentence is about what i found then talk about what the lit says then say how same/diff
- how to identify key papers 
- whatever the novelty is, find something that is most similar on the other axes. so if its a new method, find something with similar study pop/loc/yr. 


- then talk about what i did and say why we did something different.

wf synthesis
all-cause: mcbrien -reduced all-cause outpatient 1 day after PM exp but increased on 4 of 5 days after among dme users. proximity to one fire associated with reduced all-cause outpatient among dme users.
cvd: mcbrien -evac and proximity associated with increased inpt visits for 1 fire not the other among dme users. chen association between smoke events and ED visits for cvd null at earlier lags but stronger at longer lags.
resp: mcbrien -evac and proximity associated with increased inpt visits for 1 fire not the other among dme users; aguilera causal effects of \PM attributable to wildfire smoke = 10.0\% in resp admissions compared to \PM from other sources which was estimated to increase admissions by 0.76\%; chen elevated risk of ED visits for all resp diseases at lag 1, some same day too. elevaded risk of ED visits for cvd disesases at lag 10. 
mental health: cleland -wf \PM assoc w reduced same day cognitive performance; wettstein-increases in psychotropic med prescriptions in fire period; chen mixed results for mental health outcomes; jung wfs assoc with increased ED visits for all mental health conditions over lag 0-7 days.


dual exposure synthesis
do et al: dual exposure is increasing etc.
sheridan et al: ED injury effects of wind events are exacerbated with PO
deng et al: co-occurrence of PO and floods is higher than PO or floods alone for resp FWBD and CVD hosp admissions
Lin et al 2021: joint effects of storm and po are stronger than either alone. 
Lin et al 2024: impact of thunderstorms on resp disease in NYS is stronger during a PO. 
Xu et al and Yamasaki et al and lin et al: heat-related ambulance use nearly doubles during POs, esp for older adults. Stone et al finds increases in both morbidity and mortality during heat waves when there is also a power outage in three major cities during three major power outages. lin2011 also finds elevated effects from heat on resp hosp admissions and mortality during NYC 2003 blackout.



PO literature
\citep{casey2020power, northrop2024power, lin2021immediate, mcbrien2025, molinari2017s, do2025impact, zhang2020power, dominianni2018health, xiao2019immediate, ibrahim2016erratic, moreno2019community, skarha2021association, andresen2023understanding, lin2024joint, anderson2011impacts, yamasaki2024heat}

\textbf{Casey et al., Power Outages and Community Health: a Narrative Review}
\begin{itemize}
    \item power outages may affect many different health endpoints, including: carbon monoxide poisoning, temperature-related illness, GI illness, all-cause, cvd, resp, renal. 
    \item risk is heightened for those relying on electricity-dependent medical equipment, those who are older or rely on other people to complete daily activities, as well as those with underlying conditions. 
    \item most existing literature examines single, large-scale power outages. exposure was based on residence in the area of the outage, and typically involved some threshold of duration and number of customers affected. 
    \item cutpoints of people affected vary from study to study as well as depending on the season/temperature.
    \item generally, hospital visits decline right before a storm and increase during and after power outages. two studies in NY (dominianni et al and zhang et al) evaluated power outages in NY and found evidence of increases in cvd/resp/renal hospitalizations and COPD hospitalizations, respectively, following an outage.
    \item largest increases in copd hospitalizations during the first 3 days after an outage, and 23\% of copd hospitalizations on po days can be attributed to the outage. 
    \item mixed evidence of GI illness (pathway is likely po \> food refrigeration + water supply system + disinfection \> gi illness)
    \item some evidence of temp related illness both during extreme heat and extreme cold
    \item some evidence of mnch effects
    \item qualitative studies identified mental health and well being effects such as worry, anxiety,stress usually due to the disruptions caused by power outages 
    \item overall exposure assessment is usually based on residential address and has focused on single outages but has found associations between power outages and health. further research is needed as power outages usually co-occur with other disasters.
\end{itemize}

\textbf{Lin et al., The immediate effects of winter storms and power outages on multiple health outcomes and the time windows of vulnerability}
\begin{itemize}
    \item using distributed lag nonlinear models to look at winter storms/POs on hospitalizations (cvd, resp, food/water borne, injury), found joint effects of storm/PO are strongest, then storm alone, and null effects for po alone (except for injury and food/water borne being elevated on lag days 1 and 3)
    \item hospital admissions for cvd, resp, food/water borne, injury all increase during winter storm with po (RRs range from 1.04-1.90) and this is consistent with other literature. 
    \item exposure assessment: used PO data from NYS and 'estimated the coverage of po by dividing the number of customers affected over the total and then defined the po occurrence (yes/no) as the po coverage exceeding the 50th percentile of the distribution of po coverage among all operating divisions and dates with po, which was based on a prior study that detected health impact of po with similar po coverage.'
\end{itemize}

\textbf{Northrop et al. Power outages and pediatric unintentional injury hospitalizations in New York State}
\begin{itemize}
    \item exposure assignment: POL data in 30-minute increments with the number of customers served and number of customers without power for each interval. Outages were defined at the POL-level using a 36h cumulative metric, aggregating hours where $geq$10\%, $geq$20\%, $geq$50\% of customers were without power immediately preceding admission. 
    \item association between outage exposure and risk of pediatric unintentional injury hospitalization. 
    \item results: in urban non-nyc areas, outages $geq$4 hours were associated with 30\% increased odds of all-cause unintentional injury hospitalizations when $geq$50\% of customers were without power. 
    \item sensitivity analysis: using a 1\% threshold, power outages were associated with all-cause unintentional injury across all regions. This shows the impact of smaller-scale outages that may not have been captured with the predefined outage thresholds. 
\end{itemize}

\textbf{Molinari et al. Who's at risk when the power goes out? The at-home electricity-dependeng population in the United States, 2012}
\begin{itemize}
    \item In 2012, Molinari et al. estimated 366,619 people of the 175 million individuals on employer-sponsored insurance plans were reliant on electricity-dependent medical equipment (prevalence of 218.2 per 100,000). 
    \item Across all insurance types, 685,000 electricity-dependent people live at home. 
    \item Those who live at home and rely on electricity-dependent DME are particularly vulnerable. 
    \item Battery life for oxygen concentrators are typically around 3-4 hours but can be extended 8+ hours with an extra battery. 
    \item Prevalence varied by location and age, and the highest prevalence was in those 65+ and in the south and the west. 
    \item When DME fails or batteries die, individuals who live at home and rely on electricity-dependent DME account for a significant amount of avoidable ED admissions, and likely thus would account for a significant amount of morbidity/mortality associated with power outages. 
    \item 99\% of DMEs are oxygen concentrators
\end{itemize}

\textbf{Do et al., The Impact of Power Outages on Cardiovascular Hospitalizations Among Medicare Fee-for-service Enrollees in New York State, 2017--2018}
\begin{itemize}
    \item exposure assessment: using power operating localities (POL) data, they calculated the hourly average of customers without power in the POLs and then determined hourly customer averages for ZCTAs via spatial intersection. Power outage hour was defined using the percent of customers without power in a given ZCTA-hour. They used three cutpoints ($geq$1\%, $geq$10\%, and $geq$20\%). These capture a different balance of small and large outages. Exposure was binary for a given ZCTA-hour -- whether the ZCTA-hour had a PO where the percent of customers affected exceeded the cutpoint. 
    \item null results examining PO exposure and CVD hospitalizations among medicare enrollees.
    \end{itemize}

\textbf{Zhang et al. Power outage: an ignored risk factor for copd exacerbations}
\begin{itemize}
    \item copd = third leading cause of death in US
    \item examined the association between occurrence, coverage, and duration of power outages with copd hospitaliztaion
    \item po occurrence had largest association with copd hospitalizations, followed by duration then coverage. 
    \item 'this study found that the rate for overall copd hospitalization rose significantly following POs and was enhanced as PO coverage and duration were elevated.'
    \item evidence of lag effects -- impact of PO on copd hospitalizations may continue over the week after a PO. 
\end{itemize}

\textbf{Dominniani et al., Health impacts of citywide and localized power outages in New York City}
\begin{itemize}
    \item exposure assessment: the authors looked at zip codes that were affected by each outage and copmpared the number of health events that occurred during power outages with those that occurred on other days during the study period. 
    \item this study looked at power outages in NYC and used a poisson time series regression. 
    \item respiratory hospitalizations and renal disease hospitalizations were elevated following some but not all of the outages studied. all-cause mortality was associated the larger outages but not all warm-weather localized outages. cold-weather localized outages, on the other hand, were associated with elevated all casue mortality and cvd hospitalizations, but reduced resp disease hospitalizations.
\end{itemize}


\textbf{Ibrahim et al., Erratic electricity supply (Dumsor) and anxiety disorders among university students in Ghana: a cross sectional study} 
\begin{itemize}
    \item this study examined anxiety and nervousness due to the erratic energy supply in Ghana. they found that almost a quarter of students included in the study feel anxious or nervous almost every day due to the erratic power supply.
\end{itemize}

\textbf{moreno, Community resilience to power outages after disaster: A case study of the 2010 Chile earthquake and tsunami}
\begin{itemize}
    \item power outages threaten basic needs and lead to worry anxiety as a result
    \item another theme is lack of information and communication leading to uncertainty
    \item this qualitative study found that community reslience was important to mitigate the mental health impacts of power outages. this may include things like shared resources and knowledge, community kitchens, sense of community overall. 
    \item the needs that were threatened varied based on the lag after the power outage, but long term psychological disturbance can remain for longer periods. communities themselves can actively mitigate some of these negative effects, making community resilience an important piece of the equation.
\end{itemize}

\textbf{skarha, Association of power outage with mortality and hospitalizations among Florida nursing home residents after Hurricane Irma}
\begin{itemize}
    \item population: nursing home residents (~67k)
    \item exposure assessment: florida state law required a daily report from nursing homes on key data such as power outage status before during and after the hurricane, so data on whether a nursing home had power on a specific day were used as the exposure information. a dataset was created with residents living in nursing homes with power outages of any length of time during the storm, and those living in nursing homes with no power loss during that same 5 day storm period.
    \item 'power loss was associated with an increased odds of mortality within 7 and 30 days, with a roughly 25\% increase in mortality at 7 days after landfall and about a 10\% increase 30 days after landfall. Neither age nor stay status significantly modified these results. Power loss was also associated with an increased odds of first hospitalization among residents aged 65 to 74 years old.'
    \item association with power loss and hospitalization only among residents aged 65 to 74 years.
\end{itemize}

\textbf{andresen, Understanding the social impacts of power outages in North America: a systematic review}
\begin{itemize}
    \item 
\end{itemize}

\textbf{lin, The joint effects of thunderstorms and power outages on respiratory-related emergency visits and modifying and mediating factors of this relationship}
\begin{itemize}
    \item 
\end{itemize}

yamasaki2024heat
anderson2011impacts





PSPS literature
\citep{harding2024association, harding2025a, WONGPARODI2022102495, huang2023review}
\textbf{Harding et al. 2024, The association of emergency department visits and Public Safety Power Shutoffs in California}
\begin{itemize}
    \item exposure assessment: determined the numnber of customers and percent of households without power for 24 hours of each county-day. 
    \item both exposure metrics showed similar results. On days with the highest levels of PSPS event, there were increases in ED visits among older adults. additionally, on days with highest psps exposure there were more ed visits overall for respiratory, cvd, injury, and mental/behavioral diagnosis codes.
\end{itemize}

\textbf{Harding et al. 2025, The Association of Public Safety Power Shutoffs and Motor Vehicle Crashes}
\begin{itemize}
    \item PSPS events are believed to prevent a lot of fires and thus benefits are believed to outweigh costs, but few studies have been done on this 
    \item fires started by electrical infrastructure often start during high winds and thus are some of the worst fires as the winds exacerbate them 
    \item psps events thoeretically improve health and safety, but theres not really any studies on it. 
    \item psps events were not associated with increased motor vehicle crashes, despite failing traffic signals, lighting, and other safety equipment. This may be due to changes in driving habits and behavior. 
\end{itemize}

\textbf{Wong et al., Support for public safety power shutoffs in California: Wildfire-related perceived exposure and negative outcomes, prior and current health, risk appraisal and worry}
\begin{itemize}
    \item 'hazard worry was associated with increased support for PSPS, and PSPS worry was associated with decreased support for PSPS'
    \item overall, those who were more worried abotu the health impacts of wfs were more supportive of psps events, but those who were more worried about losing power were less supportive of psps. 
    \item overall, there were moderate levels of support for psps events but there is also evidence of worry and anxiety around these events. 
\end{itemize}

\textbf{Huang et al., A review of public safety power shutoffs (PSPS) for wildfire mitigation: Policies, practices, models and data sources}
\begin{itemize}
    \item 'In one way, wildfires can create electric infrastructure damage, where the risk can be mitigated by reducing the damage produced by fire events. In another way, wildfires can be caused by electric infrastructure (e.g., power lines and equipment), in which the risk can be minimized by reducing the frequency of electric infrastructure-related fires.'
    \item PSPS need to balance risks of power line ignited wfs and harms of psps 
    \item this paper has a very good summary of the history of psps events in CA and the policy measures that have led to where we are today. it is a review of the mitigation practice and how the risk mitigation tradeoff decisions are made. 
\end{itemize}


anderson2011impacts
yamasaki2024heat, xu2025power




copd literature
\citep{molinari2017s, casey2020power, zhang2020power, sutherland2005wildfire, reid2019wildfire, wilgus2024clearing}

\textbf{molinari2017s, casey2020power, zhang2020power: see above}
\textbf{Sutherland et al., Wildfire smoke and respiratory symptoms in patients with chronic obstructive pulmonary disease}
\begin{itemize}
    \item evidence that wildfire smoke (via increases in air pollution) is associated with increases in resp symptoms in patients with copd
\end{itemize}

\textbf{Reid et al., Wildfire smoke exposure under climate change: impact on respiratory health of affected communities}
\begin{itemize}
    \item current evidence shows associations between wildfire smoke exposure and respiratory health outcomes, though some of the evidence has been mixed. 
    \item the clearest association is between wfs exposure and asthma
    \item evidence of association between wfs and copd was stronger but recent studies, including some in WA and CA, have shown more mixed results. 
    \item previous mixed evidence of resp infections after wfs exposure is now stronger, though different groupings and exposure definitions make it hard to completely compare. 
    \item when all resp outcomes are grouped, there are strong relationships between wfs exposure and hospitalization
    \item there is evidence that resp health is affected by wfs exposure, but the evidence is inconsistent due to inconsistent methods in both exposure assessment, outcome definition, and statistical methods. 
\end{itemize}

\textbf{Wilgus et al., Clearing the Air: Understanding the Impact of Wildfire Smoke on Asthma and COPD}
\begin{itemize}
    \item this review walks through the pathways for wfs affecting resp health, useful if further details are needed. 
    \item wildfires are associated ed visits, hospitalizations, and worsening of symptoms for people with asthma and copd. there is also risk of long term resipiratory effects and mortality. 
    \item this review shows that evidence for asthma is strongest
    \item evidence linking wfs to cobd is less strong but there is a consensus that there is an association between wfs and copd ed or hospitalizations (though some studies do disagree)
    \item evidence of effects beyond the acute period
    \item challenges in actually quantifying the effects of wfs on resp diseases include lack of control populations, lack of long term studies, combining of resp diagnoses in many studies, inconsistent exposure assessment methods, inconsistent statistical methods
\end{itemize}




WFS literature
\citep{mcbrien_wildfire_2023, elser2025wildfire, cleland2022short, gould2024health, aguilera2021wildfire, wettstein2024psychotropic, chen2023emergency, jung2025fine}
\textbf{mcbrien_wildfire_2023, Wildfire Exposure and Health Care Use Among People Who Use Durable Medical Equipment in Southern California}
\begin{itemize}
    \item years: 2016 -- 2020
    \item study population: KPSC patients across 7 southern CA counties (kern, LA, ventura, san bernadino, orange, riverside, SD) who were 45+ years of age (as of October 28, 2019) and had rented DME equipment in the year prior. 
    \item exposure assessment: measured wildfire smoke exposure using daily wildfire and nonwildfire pm concentrations at the zcta level using tarik's ML model. then they created groups of zctas and used the average concentrations across all component zctas. to measure proximity to a wildfire, they used the wildfire boundaries and considered zctas exposed if their boundary was within 20km of a final fire perimeter on days when the fire was active. 
    \item outcome and health dataset: KPSC EHR data -- used daily counts of all-cause outpatient, ED, inpatient visits, as well as ED/inpatient visits for circulatory or respiratory disease outcomes. 
    \item statistical analysis: negative binomial regression + DiD using neg binomial regression for proximity analysis. models controlled for similar things such as temp, nonwildfire pm, ses vars. of note, the model of proximity and the outcome did not control for wildfire pm bc they did not consider this a confounder but rather part of the exposure. 
    \item major results: daily 10 micrograms per m3 increase in wildfire \PM was associated with a decrease in the rate of outpatient visits 1 day later (reduced rate [RR]=0.96) but increases on four of the five subsequent days (RR range 1.03-1.12). wildfire \PM not associated with count of all cause or inpatient visits or cardioresp ED or inpatient visits. also observed lagged effects 1-2 weeks out. For proximity results, residence in ZCTAs within 20km of the woolsey fire was associated with increased inpatient admissions for cardioresp (RR=1.45) and associated with decreased all-cause outpatient visits (RR=.89). Similar results were found for residence in a zcta located in an evac zone of the woolsey fire. Results were null for the getty fire. 
\end{itemize}

\textbf{elser2025wildfire, Wildfire smoke exposure and incident dementia}
\begin{itemize}
    \item years: January 2008 to December 2019
    \item study population: members of Kaiser Permanente Southern California (KPSC), which serves 4.7 million people across 10 California counties. KPSC members aged 60 years or older. 
    \item exposure assessment: daily mean concentrations of total \PM were estimated by CT using ML (tarik). smoky tract-days were days when the hms smoke plume boundary intersected the census tract. hmm then what? 
    \item outcome and health dataset: primary outcome was incident dementia, used EHR data from KPSC
    \item statistical analysis: discrete-time approach with pooled logistic regression. Models adjusted for age, sex, race and ethnicity (considered as a social construct rather than as a biological determinant), marital status, smoking status, calendar year, and census tract–level poverty and population density. Stratified models assessed effect measure modification by age, sex, race and ethnicity, and census tract–level poverty.
    \item major results: In adjusted models, a 1 micro-g per m3 increase in 3-year average wildfire PM2.5 concentration was not statistically significantly associated with the odds of dementia diagnosis (odds ratio [OR], 1.12; 95\% CI, 0.98-1.28). For nonwildfire PM2.5, the odds of dementia diagnosis increased by 1\% for every 1-μg/m3 increase in 3-year mean exposure (OR, 1.007; 95\% CI, 1.002-1.011).
\end{itemize}

\textbf{cleland2022short, Short-term exposure to wildfire smoke and PM 2.5 and cognitive performance in a brain-training game: a longitudinal study of US adults}
\begin{itemize}
    \item years: 2017 -- 2018 
    \item study population: The cohort included all Lumosity users in the contiguous U.S. who were ≥18 years of age and who signed up for the platform and completed between 1 and 20 plays of Lost in Migration during 1 January 2017 to 31 December 2018. The cohort included both free users and paid subscribers.
    \item exposure assessment: To estimate ambient PM2:5 for each ZIP3, we used measurements from the U.S. Environmental Protection Agency (EPA)’s Federal Reference Method/Federal Equivalent Method (FRM/FEM) monitors and the network of low-cost PurpleAir monitors. Hourly and daily average FRM/FEM observations were downloaded from the U.S. EPA’s Air Quality System database. Smoke density was determined to be none/light/medium/heavy in each zip based on smoke plume data from HMS. 
    \item outcome and health dataset: For each Lumosity user, we had data on the first three digits of the user’s ZIP code (ZIP3) location at sign-up, their self-reported age (≥18 y), gender, and education level at sign-up, as well as the device on which all games were played (iPhone, iPad, Android, or web). For each gameplay, data included the user’s score at the end of the game (raw score) and their score relative to a normative population of Lumosity users sampled to match the 2010 U.S. Census for age, gender, and education level (percentile score). We used the raw score (hereafter referred to as the attention score) as the primary measure of cognitive performance and considered the percentile score in a sensitivity analysis.
    \item statistical analysis: unconstrained distributed lag model with 7 lags of pop weighted daily avg pm, where lag 0 is the concentration on the day of gameplay, lag 1 is the day prior, etc. for subdaily pm exposure, used the max pop weighted hourly avg pm concentration in the 3/6/12h prior to game play. smoke exposure looked at the same lags but could not do subdaily due to availability. 
    \item major results: Daily and subdaily \PM were negatively associated with attention score. A 10 micro-g per m3 increase in \PM in the 3 h prior to gameplay was associated with a 21.0 point decrease in score. \PM exposure over 20 plays accounted for an estimated average 3.7\% reduction in final score. Associations were more pronounced in the wildfire-impacted western U.S. Medium and heavy smoke density were also negatively associated with score. Heavy smoke density the day prior to gameplay was associated with a 117.0 point decrease in score relative to no smoke.
\end{itemize}

\textbf{gould2024health, Health effects of wildfire smoke exposure}
\begin{itemize}
    \item meta analysis of wf and health studies
    \item "While accumulating evidence makes it clear that inhaled wildfire smoke negatively impacts human health, it is also increasingly clear that wildfire smoke is different from pollution from other sources in ways that likely matter for human health."
    \item Same-day all-cause mortality increased by 0.15\% [95\% confidence interval (CI) 0.01–0.28\%] per 1–mg m 3 increase in wildfire specific PM2.5. There were robust positive associations between wildfire PM2.5 and same-day respiratory outcomes: Respiratory hospitalizations increased by 0.25\% (95\% CI 0.09–0.52\%) and respiratory ED visits increased by 0.36\% (95\% CI 0.19–0.53\%) per additional 1–mg m 3 increase in ambient wildfire smoke PM2.5. We found a non–statistically significant 0.06\% (95\% CI 0.00–0.12\%) increase in same-day cardiovascular hospitalizations and no meaningful change in same-day cardiovascular ED visits ( 0.03\%; 95\% CI 0.18–0.12\%) per additional 1–mg m 3 increase in ambient wildfire smoke PM2.5. For all outcomes except respiratory hospitalizations and cardiovascular ED visits, there was evidence of heterogeneity in effects across studies (i.e., Q-statistic p < 0.05). Egger’s tests did not indicate evidence of publication bias for any outcome (p > 0.05)
    \item "Future work should aim to leverage and understand the unique dynamics of wildfire smoke: (a) Temporal variation in ambient wildfire smoke concentrations is frequently idiosyncratic, enabling causal interpretations when appropriately controlling for area characteristics and time trends; (b) healthcare utilization may respond nonlinearly to increasing pollution levels due to a combination of individual-level pathophysiological impacts and changing behaviors at the individual and population level; and (c) healthcare utilization and health impacts."
\end{itemize}

\textbf{aguilera2021wildfire, Wildfire smoke impacts respiratory health more than fine particles from other sources: observational evidence from Southern California}
\begin{itemize}
    \item studies suggest that pm from wildfires is more toxic than pm from other sources: "All the above compounds in wildfire smoke tend to generate more free radicals and thus have a greater potential to cause inflammation and oxidative stress in the lung than urban ambient particulate from the same region"
    \item In order to robustly isolate the health impacts of wildfire-specificPM2.5, we apply and compare four analytical approaches. Specifically, we implemented the following distinct methods: (i) an instrumental variable approach with a two-stage regression; (ii) a spatio-temporal multiple imputation approach, and (iii) an interaction effect approach. Lastly, we also compared these approaches to (iv) a seasonal interpolation method recently proposed by Lipner et al.
    \item years: 1999-2012
    \item study population: Southern California
    \item exposure assessment: they used \PM data from US EPA AQS at ground monitoring stations within 20km of a pop weighted zipcode centroid. wildfire exposure came from CalFire. And smoke plumes came from NOAA HMS
    \item outcome and health dataset: Daily hospital admissions for respiratory diseases were obtained from the California Office of Statewide Health Planning and Development (OSHPD)
    \item statistical analysis: they used 4 different ways to estimate wildfire specific \PM (see bullet above). 
    \item major results: Our findings indicate that wildfire-specificPM2.5 can cause a greater impact on respiratory health than \PM from other sources.
    \item `Based on the mean number (1.85) of daily respiratory admissions per 100,000 individuals, a 10 micro-g per m3 increase in PM2.5 was estimated to increase the number of admissions by only 0.76\%(95\%CI:0.42–1.1). In contrast, the causal effects of \PM attributable to wildfire smoke estimated by spatio-temporal imputation amounted to a 10.0\% (95\% CI: 3.5–16.5) increase in admissions, the highest percentage among all methods and exposures used. Such results were similar, though varying in the amplitude of \% increase in admissions when using other approaches to isolate the wildfire-specific \PM, as well as among all approaches considering smoke plumes within a 160 km buffer to define zip code days exposed to wildfire smoke. We conclude that wildfire-specificPM2.5 is up to 10 times more harmful on human health than PM2.5 from other sources.'
\end{itemize}

\textbf{wettstein2024psychotropic, Psychotropic medication prescriptions and large California wildfires}
\begin{itemize}
    \item years: 2011 -- 2018
    \item study population: people residing in metropolitan statistical areas with prescriptions for psychotropic meds
    \item exposure assessment: proximity to large wildfires in California using data from CalFire redbook
    \item outcome and health dataset: prescriptions of psychotropic medications, including antidepressants, antipsychotics, anxiolytics, hypnotics, and mood-stabilizers, with statins as a negative control outcome. Data came from the Merative MarketScan database. Used daily medication prescriptions from marketscan and information about the met stat area that each patient resided in in addition to individual level characteristics. 
    \item statistical analysis: cohort study used an interrupted time-series analysis to compare psychotropic medication prescriptions before/after wildfires in CA. they categorized prescriptions into 7 different medication groups.  Using an ITS regression framework they used the first twitter posting date for each fire as the point at which it started and then did their ITS for 6 weeks before/after. Used poisson generalized linear models 
    \item major results: increase in mean daily medication prescriptions for the fire period, for antidepressant medications (1.04 [95\% CI, 1.01-1.07]), anxiolytic medications (1.05 [95\% CI, 1.02-1.09]), mood-stabilizing medications (1.06 [95\% CI, 1.01-1.13]), and all psychotropic medications (1.04 [95\% CI, 1.01-1.07])
    \item psychotropic medication prescriptions increased in the fire period compared with the prefire baseline, including for antidepressant medications, anxiolytic medications, and mood-stabilizer medication, but not for antipsychotic medications, hypnotic medications, or statin medications, the nonpsychotropic control.
\end{itemize}

\textbf{chen2023emergency, Emergency department visits associated with wildfire smoke events in California, 2016--2019}
\begin{itemize}
    \item investigated associations between short-term exposure to smoke events, derived from wf specific pm, and EDVs for resp, cardio, metabolic, and mental health outcomes 
    \item years: 2016 -- 2019, restricted analysis to June -- December each yr. 
    \item study population: California 
    \item exposure assessment: wildfire \PM came from tarik's group and smoke plume data from NOAA's HMS identified smoke-exposed days. same method as elser 2023 to get wf pm and non wf pm. restricted analysis to days with high concentrations of wf-specific \PM. to do so, defined a smoke event as a day with a wf \PM concentration at or above the 98th percentile across all days during June-Dec 2016-2019.
    \item outcome and health dataset: used daily counts of unscheduled EDVs from California Department of Health Care Access and Information (HCAI). 
    \item statistical analysis: two stage time series analysis using quasi-poisson regression with lag effects and random effects
    \item major results: Smoke events were associated with an increased risk of EDVs for all respiratory diseases at lag 1 [14.4\%, 95\% confidence interval (CI): (6.8, 22.5)], asthma at lag 0 [57.1\% (44.5, 70.8)], and chronic lower respiratory disease at lag 0 [12.7\% (6.2, 19.6)]. Smoke events were also associated with a 3.2\% (95\% CI: 1.7, 4.7\%) increase in the risk of EDVs for all cardiovascular diseases at lag 10. The association between smoke events and EDVs for all cardiovascular diseases was negative or null during earlier lags but gradually became positive with longer lags, suggesting that cardiovascular impacts may be delayed compared to respiratory impacts. Mixed results were observed for mental health outcomes. There were mixed results for mental health outcomes. 
    \item "Our observation of strong, immediate respiratory effects and more delayed cardiovascular effects may be partly explained by the timing of the physiological effects of wildfire smoke."
\end{itemize}

\textbf{jung2025fine, Fine particulate matter from 2020 California wildfires and mental health--related emergency department visits}
    \item years: July 2020 -- December 2020
    \item study population: California residents
    \item exposure assessment: wildfire \PM from Childs et al. 
    \item outcome and health dataset: used daily counts of unscheduled EDVs from California Department of Health Care Access and Information (HCAI). 
    \item statistical analysis: distributed lag nonlinear model for lags 0-7. 
    \item major results: A 10-μg/m3 increase in wildfire-specific PM2.5 was associated with higher ED visits for all-cause mental conditions (cumulative relative risk [cRR] over lag 0-7 days, 1.08; 95\% CI, 1.03-1.12), depression (cRR over lag 0-7 days, 1.15; 95\% CI, 1.02-1.30), other mood-affective disorders (cRR over lag 0-7 days, 1.29; 95\% CI, 1.09-1.54), and anxiety (cRR over lag 0-4 days, 1.06; 95\% CI, 1.00-1.12). EMM by race. 
    \item Wildfire smoke exposure was associated with significantly increased odds of subsequent ED visits for mental health conditions in this cross-sectional study, with varying lag times for different subconditions and demographic groups.




Dual exposure literature
\citep{do2025spatiotemporal, sheridan2021individual, deng2022independent, lin2021immediate, lin2024joint, xu2025power, yamasaki2024heat, stone2023blackouts}

\textbf{do2025spatiotemporal, Spatiotemporal patterns of individual and multiple simultaneous severe weather events co-occurring with power outages in the United States, 2018--2020}
    \item years: 2018 -- 2020
    \item study population: N/A
    \item exposure assessment: For wildfires, we used data from the National Interagency Fire Center, and a county-day was considered exposed when it overlapped with a \geq 1km2 actively burning wildfire at any point of the day. For power outages, used poweroutage.us data. aggregated data to the county hourly level as the avg count of customers without power and included counties with more than 50\% coverage of county wide electrical customers and 50\% reporting for the full 3 year period -- defined as counties with reliable data. calculated customer coverage as the percent of eleictrical customers tracked by power outage.us for each county, which was the total electrical customers recorded in the poweroutage.us data over the total electrical customers estimated using data from the energy information administration 861 forms. allocated total customers served in a state to counties based on proportion of households and businesses in each county reported by the census. calcualted percent of hours tracked in each county. looked at outages lasting 8+ hours and defined this as when the percent of customers without power in a county hour exceeded 0.1\% for at least 8 consecutive hours. to calcualte the percent of customers without power for every county hour, divided the customers without power reported from poweroutage.us by the total customers served in that county based on the 861 forms. 
    \item outcome and health dataset: N/A
    \item statistical analysis: analysis was at the countty-day level so they identified county days with 8+ hour power outage that happened on same dat as individual or multiple severe weather events. described nationwide spatiotemporal patterns ofindividual (e.g., tropical cyclone alone) and simultaneous (e.g., tropical cyclone +anomalous heat) severe weather events co-occurring with 8+ hour outages. We used hourly county-level PowerOutage.us data from 2018–2020 todefine 8+ hour outages as whenever the daily proportion ofcustomers without power was \geq 0.1\% for \geq 8 continuous hours. We conducted analyses at the daily and county (county-day) level and identified county-days with severe weather events, including anomalous cold, anomalous heat, anomalous precipitation, snowfall, tropical cyclones, and wildfire.
    \item major results: Of 1,657 counties with reliable outage data, 1,205 (72.7\%) experienced an 8+ hour outage co-occurring with an individual severe weather event, and 904 (54.6\%) with multiple simultaneous severe weather events. On the West Coast, wildfires were increasingly likely toco-occur with weather-related outages from 2018–2020. Over the 3-year period, we noted asuggestive upward trend in the proportion of wildfires co-occurring with power outages relative to other severe weather events in California (2018: 24\%, 2019: 29\%, 2020: 28\%), Oregon (2018: 10\%, 2019: 0\%, 2020: 34\%), and Washington (2018: 9\%, 2019: 0\%, 2020: 14\%). Over the 3-year period, we noted asuggestive upward trend in the proportion of wildfires co-occurring with power outages relative to other severe weather events in California (2018: 24\%, 2019: 29\%, 2020: 28\%), Oregon (2018: 10\%, 2019: 0\%, 2020: 34\%), and Washington (2018: 9\%, 2019: 0\%, 2020: 14\%).

\textbf{sheridan2021individual, The individual and synergistic impacts of windstorms and power outages on injury ED visits in New York State}
    \item years: November–April 2005–2013
    \item study population: NYS
    \item exposure assessment: Wind event data was obtained from NOAA severe weather inventory, and was defined according to high wind, strong wind or thunderstorm wind defined by NOAA. Used NYS department of public service PO data. These data include the total number of customers in each power operating dividion. then computed the PO coverage as the percent of customers affected. Then among all the records (i.e. division*day) with PO coverage >0, we chose the 50th percentile of PO coverage as the threshold. We defined PO occurrence for each record (i.e. division*day) based on whether the PO coverage exceeded the threshold.
    \item outcome and health dataset: SPARCS: emergency department (ED) visit data were obtained from the NYS Department of Health from  to identify injury cases, length of stay and care costs. used all emergency department (ED) visits, both inpatient and outpatient
    \item statistical analysis: By comparing non-event days, we used distributed lag nonlinear models to evaluate the impacts of wind events, PO, and their combined effect on injuries during the cold season over a 0-3-day lag period, while controlling for time-varying confounders. The differences in critical care indicators between event and non-event days were also evaluated.
    \item major results: Overall injuries ED visits (16,628,812) significantly increased during the wind events (highest Risk Ratio (RR): 1.05; 95\% CI: 1.02–1.08), and were highest when wind events cooccurred with PO (highest RR: 1.14; 95\% CI: 1.10–1.18), but not during PO alone (RR: 1.00; 95\%CI: 0.96–1.04). The increase was also observed with all subgroups through Day 2 after the event. Greater risks exist for older adults (≥65 years) and those on Medicaid. After the joint occurrences of wind events and PO, average visits are 0.2 days longer, and cost 13\% more, compared to no wind/no PO days.

\textbf{deng2022independent, The independent and synergistic impacts of power outages and floods on hospital admissions for multiple diseases}
    \item years: 2002 -- 2018
    \item study population: New York State
    \item exposure assessment: power outage data from NYS department of public service, flood events from NOAA (storm events database) which includes storms and other weather phenomena with enough intensity to cause significant damage. Assigned each weather event to the corresponding power-operating division. used PO data from the NYS department of public service, including total number of customers in each power operating division, the dates of PO, and the numbers of customers affected by PO. We estimated the coverage of PO by dividing the number of customers affected over the total and then defined the PO occurrence (Yes/No) as the PO coverage exceeding the 75th percentile of the distribution of PO coverage (1.72\%) among all operating divisions and dates with PO. This definition of PO occurrence was based on previous studies (Sheridan et al., 2021; Lin et al., 2021; Zhang et al., 2020). Overall, health data, event data, and PO data were linked at the power operating division level.
    \item outcome and health dataset: SPARCS data on hospital admissions, used ICD codes to identify their otucomes of interest (cvd, resp, food and water borne diseases). The health outcome was the sum of the daily number of hospital admissions for each disease at the operating division level.
    \item statistical analysis: distributed lag nonlinear models
    \item major results: significant increased health risks associated with both the independent effects from PO and floods, and their synergistic effects. Generally, the Rate Ratios (RRs) for the co-occurrence of PO and floods were the highest, followed by PO alone, and then floods alone, especially when PO coverage is >75th percentile of its distribution (1.72\% PO coverage). For PO and floods combined, immediate effects (lag 0) were observed for chronic respiratory diseases (RR:1.58, 95\% CI: 1.24, 2.00) and FWBD (RR:3.02, 95\%CI: 1.60, 5.69), but delayed effects were found for cardiovascular diseases (lag 3, RR:1.13, 95\%CI: 1.03, 1.24) and respiratory infections (lag 6, RR:1.85, 95\%CI: 1.35, 2.53). The risk association was slightly stronger among females, whites, older adults, and uninsured people but not statistically significant.

\textbf{lin2021immediate}
see above.

\textbf{lin2024joint}
see above.

\textbf{xu2025power, Power Outages Increase the Risk of Ambulance Attendances Associated with Non-Optimal Temperatures}
    \item years: 2012 -- 2016
    \item study population: Victoria, Austrailia. Split up older and younger than 65 but no specific info on age is provided. 
    \item exposure assessment: Used daily mean temp and humidity from ERA5. Used daily nighttime light products to estimate the power loss across all communities. after processing the nighttime light data, they used a robust linear regression to do XX. They then calculated power supply and power loss metrics using the nighttime lights data for use in their model.  
    \item outcome and health dataset: Daily emergency ambulance attendances from 79 LGAs (regions) were obtained from Ambulance Victoria, which is the state wide emergency medial service in Victoria. These data had ICD codes and some demographic info. 
    \item statistical analysis: assessed the impacts of non-optimal temperatures on all-cause and cause-specific ambulance attendances using a case time series design and quasi-poisson regression. used an interaction term between temp/PO in the model. 
    \item major results: cold-related associations remained stable, we observed elevated heat-related risks for all-cause ambulance attendances during power outages. heat-related risk nearly doubled during power outages, leading to an additional 1110 ambulance cases and accounting for approximately 3.75\% (95\% CI: 0.48\%, 6.9\%) of heatrelated ambulance attendances. Subgroup analysis revealed a significant increase in heat-related risks during extreme heat among individuals aged 65 and older, females, and those presenting with abnormal symptoms.

\textbf{yamasaki2024heat, Heat-related mortality and ambulance transport after a power outage in the Tokyo metropolitan area}
    \item years: September 2019 typhoon (study per: july - sept 2019)
    \item study population: seven prefectures in the Tokyo metropolitan areas
    \item exposure assessment: heat related illness assessed through icd codes; power outage exposure assessed using tokyo electric power company power grid data. They calculated electricity reduction (ER) caused by the power outages at the fire department district level. each power outage case at the household level was used as a binary indicator of the existence of a power outage. ER was defined as the magnitude of a power outage when power outage cases (each case with more than 12 hours of a power outage) were >0; we assumed that a power outage event ended when the number of power outage cases reached zero. We collected daily weather data (e.g., daily mean temperature and relative humidity) from the Japan Meteorological Agency. The daily mean temperature in 2019 is comparable across the recent of 2015–2019.
    \item outcome and health dataset: collected daily caounts of heat illness from the Ministry of Health, Labour and Welfare of Japan. 
    \item statistical analysis: role of temperature during the power outage, we then examined whether the temperature–HIAT and mortality associations were modified by different power outage levels (0\%, 10\%, and 20\% ER). We computed the ratios of relative risks to compare the risks associated with various ER values to the risks associated without ER.
    \item major results: Event study results showed that the incidence rate ratio was 2.01 (95\% confidence interval [CI] = 1.42, 2.84) with effects enduring up to 6 days, and 1.11 (95\% CI = 1.02, 1.22) for mortality on the first 3 days after the typhoon hit. Comparing 20\% to 0\% ER, the ratios of relative risks of heat exposure were 2.32 (95\% CI = 1.41, 3.82) for HIAT and 0.95 (95\% CI = 0.75, 1.22) for mortality. 20\% ER was associated with a two-fold greater risk of HIAT because of summer heat during the power outage, but there was little evidence for the association with all-cause mortality.

\textbf{stone2023blackouts, How blackouts during heat waves amplify mortality and morbidity risk}
    \item years: The resulting heat wave periods identified were August 14−18, 1995 (Atlanta); June 15−19, 1994 (Detroit); and July20−24, 2006 (Phoenix).
    \item study population: Atlanta, Detroit, Pheonix 
    \item exposure assessment: derived individually experienced termpature and then looked at concurrent heat waves + blackouts and the probability of heat-related mortality or morbidity for each resident during the 5 day heat wave. Calculated heat-related attributable burdens from the baseline health burdens and then looked at the association between extreme heat days and health event. 
    \item outcome and health dataset: category-specific incidence rate obtained for each health outcome: mean daily all-natural-cause mortality and cause-specific ED visits.
    \item statistical analysis: Exposure response functions to capture effects of blackout conditions on heat risk, followed by calculation of total daily burdens and derivation of attributable fractions. 
    \item major results: substantial increase in heat-related mortality and morbidity across each of the three cities with a loss of electrical power during heat wave conditions. The estimated rate of heat-related mortality (per 100,000 population) for each city during simulated heat wave conditions with an operational electrical grid (Power On scenario) is 0.5 in Atlanta (total =2), in Detroit (total = 107), and 1.3 in Phoenix (total =19). 
